\documentclass[aspectratio=169]{beamer}

\usepackage[utf8]{inputenc}
\usepackage[brazil]{babel}
\usepackage{graphicx}
\usepackage{hyperref}

% Tema simples e limpo
\usetheme{Madrid}
\usecolortheme{default}
\setbeamertemplate{navigation symbols}{}

% -------------------------------------------------
% INFORMAÇÕES DO EP (OS ALUNOS DEVEM EDITAR)
% -------------------------------------------------

\title{MAC0209 --- Modelagem e Simulação}
\subtitle{Exercício-Programa (EP) --- Título do EP}
\author{Grupo G? \\[4pt]
Aluno 1 \\ 
Aluno 2 \\ 
Aluno 3}
\institute{Bacharelado em Ciência da Computação \\ IME -- USP}
\date{\today}

% -------------------------------------------------

\begin{document}

% =================================================
% SLIDE 1
% =================================================
\begin{frame}
\titlepage
\end{frame}

\begin{frame}{Contexto e Objetivo do EP}

\textbf{Contexto do curso}

\begin{itemize}
    \item Assunto tratado na última aula.
\end{itemize}

\vspace{0.5cm}

\textbf{Objetivo do EP}

\begin{itemize}
    \item Descrever aqui qual problema o exercício resolve.
\end{itemize}

\end{frame}

% =================================================
% SLIDE 2
% =================================================
\begin{frame}{Trabalho Realizado}

\textbf{Conceitos e implementação}

\begin{itemize}
    \item Descrição dos conceitos estudados.
    \item Descrição do experimento realizado.
\end{itemize}

\end{frame}

% =================================================
% SLIDE 3
% =================================================
\begin{frame}{Resultados e Conclusões}

\textbf{Resultados}

\begin{itemize}
    \item Principais resultados obtidos.
    \item Interpretação do comportamento observado.
\end{itemize}

\end{frame}

\end{document}
